\section{What Is New: Four Findings}


The rates of Table~\ref{tab:III} say that the policies are unsafe; the ablations say \emph{how}, and it is there that execution-phase safety departs from collision avoidance.

\textbf{(i) A safety command changes whether the task gets done, not how.} Replacing the neutral instruction with an explicit safety command leaves the violation rate unchanged at \emph{N} = 20--24 paired seeds (T1 33 \% $\rightarrow$ 29 \%, McNemar \emph{p} = 1.0; T3 33 \% $\rightarrow$ 46 \%, \emph{p} = 0.58; T6 30 \% $\rightarrow$ 25 \%, \emph{p} = 1.0; Fig. \ref{fig:fixability}). On the path this probe is ceiling-limited, so we calibrated a placement with headroom --- the stove 0.28 m off the path, blind rate 37 \% (11/30, two seeds) --- and ran the naming $\times$ rendering design (Table~\ref{tab:XI}): naming leaves the rate at 29 \% (6/21, Fisher \emph{p} = 0.76), and at 16 \% against 21 \% with the stove hidden (\emph{p} = 0.77; McNemar on 19 pairs \emph{p} = 1.0). The command reaches the policy: named-plus-visible halves completion (29 \% against 47--68 \% in the other arms, \emph{p} $\leq$ 0.001), and on $\pi_{0.5}$ a spatial command leaves the plow-through at 88 \% vs 94 \% while cutting success to 62 \%, a keep-upright command leaves its tilt as it was (16/24), and a blades-away command leaves the presentation unchanged (20/20).

\textbf{(ii) A visible hazard pulls the path toward it.} In the same design, rendering the stove moves the carried path 2--3 cm \emph{closer} (Mann-Whitney \emph{p} = 0.013 blind, 0.005 named; 33 \% vs 18 \% violating, pooled over naming), and on $\pi_{0.5}$ a rendered marker is crossed as often as an unrendered point (16/16), and one moved 0.28 m off the transport, clear of a direct carry, is entered on 12/32. Perception reaches the path, as attraction rather than avoidance.

\textbf{(iii) Neither orientation nor speed is conditioned on the person.} The carry yaw is the same at every bystander azimuth for GR00T and on either side of the table for $\pi_{0.5}$ carrying scissors (T3), the same whether the person is a capsule, a photorealistic human or not rendered (Appendix E.8), and the carry speed is the same with and without the person for both (T5a): orientation, which no stop can correct, and speed, which a slowdown would change, are never adapted. It follows the object's initial pose instead: turning the scissors 180$^{\circ}$ moves the violation to the other side.

\textbf{(iv) A moving person is walked into, and pressed against once they stop.} No deceleration precedes contact at any crossing speed (T6), and a person who stops on contact is treated as an obstacle: the payload stays pressed against them, as a coworker's hand in the bowl is pressed by $\pi_{0.5}$'s mug. The external stop that prevents it is itself incomplete on a humanoid (Appendix E.7).

\textbf{Hypothesis.} These findings are what one expects if execution-phase competences are \emph{absent from the imitation training distribution} (\S2) rather than from the prompt or the percept; the argument, its first test and four alternative readings are in Appendix F.

